%---------------------------------------------------------------------
% Part two: Surah Aal-e-Imran. Six questions.
%---------------------------------------------------------------------

\section*{حصہ دوم --- سورۃ آل عمران}
\markboth{سورۃ آل عمران}{}

\begin{nukta}[سورت کا خاکہ]
\textbf{مدنی} سورت، قرآن کی \textbf{تیسری} سورت، \textbf{دو سو} آیات، بیس
رکوع۔ نصفِ اول اہلِ کتاب سے مکالمہ اور نجران کے وفد کے واقعے پر، نصفِ ثانی
غزوہ احد کے تفصیلی تجزیے پر مشتمل ہے (آیات \textenglish{121--179})۔
\end{nukta}

%---------------------------------------------------------------------
\begin{sawal}{سوال \textenglish{10} \ \textnormal{\small(خزاں \textenglish{2021}، سوال \textenglish{5})}}
محکمات و متشابہات کے تصور پر آیت کا ترجمہ و تفسیر لکھیے۔
\end{sawal}

\begin{hal}
\ayat{هُوَ الَّذِي أَنزَلَ عَلَيْكَ الْكِتَابَ مِنْهُ آيَاتٌ مُّحْكَمَاتٌ
هُنَّ أُمُّ الْكِتَابِ وَأُخَرُ مُتَشَابِهَاتٌ فَأَمَّا الَّذِينَ فِي
قُلُوبِهِمْ زَيْغٌ فَيَتَّبِعُونَ مَا تَشَابَهَ مِنْهُ ابْتِغَاءَ الْفِتْنَةِ
وَابْتِغَاءَ تَأْوِيلِهِ وَمَا يَعْلَمُ تَأْوِيلَهُ إِلَّا اللَّهُ
وَالرَّاسِخُونَ فِي الْعِلْمِ يَقُولُونَ آمَنَّا بِهِ}{آل عمران،
\textenglish{7}}
\tarjuma{وہی ہے جس نے تجھ پر کتاب اتاری، اس میں کچھ آیات محکم (واضح المعنی)
ہیں جو کتاب کی اصل ہیں، اور کچھ متشابہ (مشتبہ) ہیں۔ جن کے دلوں میں کجی ہے وہ
فتنے اور اپنی تاویل کی طلب میں متشابہ کے پیچھے پڑتے ہیں، حالانکہ اس کی حقیقی
تاویل اللہ کے سوا کوئی نہیں جانتا، اور علم میں پختہ لوگ کہتے ہیں: ہم اس پر
ایمان لائے۔}

\medskip
\textbf{تشریح:} \textbf{محکمات} وہ آیات جن کا مفہوم واضح ہے اور جن پر عقیدہ و
عمل کی بنیاد رکھی جاتی ہے۔ \textbf{متشابہات} وہ جن کا حقیقی مفہوم انسانی عقل
سے پرے ہے (جیسے روحِ حیات کی حقیقت، یا اللہ کی صفاتِ متشابہہ کا کیف)۔ آیت دو
رویوں کا موازنہ کرتی ہے: \emph{زیغ} (کجی) والے متشابہ کے پیچھے پڑ کر فتنہ
پھیلاتے ہیں؛ \emph{راسخون فی العلم} (پختہ علم والے) اس پر ایمان لا کر معاملہ
اللہ کے سپرد کر دیتے ہیں۔ یہی صحیح طرزِ عمل ہے۔
\end{hal}

%---------------------------------------------------------------------
\begin{sawal}{سوال \textenglish{11} \ \textnormal{\small(خزاں \textenglish{2021}، سوال \textenglish{4})}}
وحدتِ امت اور تفرقے کی ممانعت میں آیت کا ترجمہ و تفسیر لکھیے۔
\end{sawal}

\begin{hal}
\ayat{وَاعْتَصِمُوا بِحَبْلِ اللَّهِ جَمِيعًا وَلَا تَفَرَّقُوا وَاذْكُرُوا
نِعْمَتَ اللَّهِ عَلَيْكُمْ إِذْ كُنتُمْ أَعْدَاءً فَأَلَّفَ بَيْنَ
قُلُوبِكُمْ فَأَصْبَحْتُم بِنِعْمَتِهِ إِخْوَانًا}{آل عمران، \textenglish{103}}
\tarjuma{اور سب مل کر اللہ کی رسی کو مضبوطی سے تھام لو اور تفرقے میں نہ پڑو،
اور اپنے اوپر اللہ کے اس احسان کو یاد کرو جب تم ایک دوسرے کے دشمن تھے، پھر اس
نے تمہارے دلوں کو جوڑ دیا اور تم اس کے فضل سے بھائی بھائی بن گئے۔}

\medskip
\textbf{تشریح:} \emph{حبل اللہ} (اللہ کی رسی) سے مراد قرآن اور دینِ اسلام ہے۔
لفظ \ar{جَمِيعًا} پر غور کیجیے --- حکم اجتماعی ہے، انفرادی نہیں۔ دلیل تاریخی
دی گئی: اوس و خزرج قبیلے صدیوں سے برسرِ پیکار تھے، اسلام نے انہیں بھائی بنا
دیا۔ آگے آیت \textenglish{105} تفرقے کی مذمت میں اہلِ کتاب کی مثال دیتی ہے:
\ar{وَلَا تَكُونُوا كَالَّذِينَ تَفَرَّقُوا وَاخْتَلَفُوا مِن بَعْدِ مَا
جَاءَهُمُ الْبَيِّنَاتُ} --- ان جیسے نہ ہو جاؤ جو واضح دلائل آ جانے کے بعد
تفرقے میں پڑ گئے۔
\end{hal}

%---------------------------------------------------------------------
\begin{sawal}{سوال \textenglish{12} \ \textnormal{\small(بہار \textenglish{2013}، سوال \textenglish{7})}}
دعوت و تبلیغ کی ضرورت و اہمیت پر آیت کا ترجمہ اور تفصیلی نوٹ لکھیے۔
\end{sawal}

\begin{hal}
\ayat{وَلْتَكُن مِّنكُمْ أُمَّةٌ يَدْعُونَ إِلَى الْخَيْرِ وَيَأْمُرُونَ
بِالْمَعْرُوفِ وَيَنْهَوْنَ عَنِ الْمُنكَرِ وَأُولَٰئِكَ هُمُ
الْمُفْلِحُونَ}{آل عمران، \textenglish{104}}
\tarjuma{اور تم میں سے ایک جماعت ایسی ہونی چاہیے جو بھلائی کی طرف بلائے،
نیکی کا حکم دے اور برائی سے روکے۔ اور یہی لوگ فلاح پانے والے ہیں۔}

\medskip
\textbf{تشریح:} آیت میں تین الگ ذمہ داریاں گنوائی گئیں: \textbf{دعوتِ خیر}
(بھلائی کی طرف بلانا)، \textbf{امر بالمعروف} (نیکی کا حکم دینا)، اور
\textbf{نہی عن المنکر} (برائی سے روکنا)۔ لفظ \ar{مِّنكُمْ} پر غور کیجیے ---
یہ کام \emph{پوری امت} پر بحیثیتِ فرضِ کفایہ ہے: ایک باقاعدہ جماعت اسے
سرانجام دے تو باقی امت بری الذمہ ہو جاتی ہے، مگر جماعت کا وجود ضروری ہے۔
فلاح (کامیابی) کو انہی تین کاموں سے مشروط کیا گیا --- محض عبادت کافی نہیں،
معاشرتی اصلاح بھی لازم ہے۔ حدیث میں اسی کی مزید تاکید ہے: ''جو برائی دیکھے،
ہاتھ سے روکے؛ نہ ہو سکے تو زبان سے؛ نہ ہو سکے تو دل سے (بغض رکھے) --- اور یہ
کمزور ترین ایمان ہے'' (مسلم)۔
\end{hal}

%---------------------------------------------------------------------
\begin{sawal}{سوال \textenglish{13} \ \textnormal{\small(بہار \textenglish{2013}، سوال \textenglish{6})}}
سورۃ آل عمران کی روشنی میں غزوہ احد کے واقعے کی تفصیل بیان کریں۔
\end{sawal}

\begin{hal}
سورۃ آل عمران کی آیات \textenglish{121} سے \textenglish{179} تک، یعنی
تقریباً ایک تہائی سورت، اسی غزوے کا تجزیہ ہے --- شوال \textenglish{3}ھ، جبلِ
احد۔ قریش تین ہزار (قیادت ابو سفیان)، مسلمان ایک ہزار --- تین سو عبد اللہ بن
ابی کے ساتھ راستے سے لوٹ گئے، میدان میں صرف \textbf{سات سو} رہ گئے۔

\medskip
\textbf{واقعہ۔} نبی کریم \saw نے پچاس تیر اندازوں کو ایک درّے پر مقرر کیا،
حکم دو ٹوک تھا: \emph{فتح ہو یا شکست، جگہ نہ چھوڑنا}۔ ابتدا میں مسلمانوں کو
فتح ہوئی اور وہ مالِ غنیمت تک پہنچ گئے؛ اکثر تیر انداز یہ سمجھ کر کہ جنگ ختم
ہو گئی، درّہ چھوڑ کر غنیمت سمیٹنے لگے۔ خالد بن ولید (جو اس وقت مسلمان نہیں
ہوئے تھے) نے خالی درّے سے پیچھے آ کر حملہ کیا اور مسلمان دو پاٹوں میں آ گئے۔
ستر صحابہ شہید ہوئے، جن میں حضرت حمزہؓ بھی تھے؛ نبی کریم \saw خود زخمی ہوئے۔

\medskip
\textbf{قرآنی تجزیہ۔}
\ayat{وَلَقَدْ صَدَقَكُمُ اللَّهُ وَعْدَهُ إِذْ تَحُسُّونَهُم بِإِذْنِهِ
حَتَّىٰ إِذَا فَشِلْتُمْ وَتَنَازَعْتُمْ فِي الْأَمْرِ وَعَصَيْتُم مِّن
بَعْدِ مَا أَرَاكُم مَّا تُحِبُّونَ}{آل عمران، \textenglish{152}}
\tarjuma{اور اللہ نے تم سے اپنا وعدہ سچا کر دکھایا جب تم اس کے حکم سے انہیں
کاٹ رہے تھے، یہاں تک کہ جب تم نے کمزوری دکھائی، حکم میں جھگڑا کیا اور نافرمانی
کی --- اس کے بعد کہ اللہ نے تمہیں وہ دکھا دیا جو تم چاہتے تھے۔}

\smallskip
تین اسباب صاف لکھے ہیں: \textbf{فشل} (کمزوری)، \textbf{تنازع} (جھگڑا)،
\textbf{عصیان} (نافرمانی)۔ حوصلہ بھی دیا گیا: \ar{وَلَا تَهِنُوا وَلَا
تَحْزَنُوا وَأَنتُمُ الْأَعْلَوْنَ إِن كُنتُم مُّؤْمِنِينَ}
(\textenglish{139})، اور شہداء کے بارے میں: \ar{بَلْ أَحْيَاءٌ عِندَ
رَبِّهِمْ يُرْزَقُونَ} (\textenglish{169})۔ \textbf{سبق:} حکم کی خلاف ورزی
فتح کو شکست میں بدل دیتی ہے۔
\end{hal}

%---------------------------------------------------------------------
\begin{sawal}{سوال \textenglish{14} \ \textnormal{\small(خزاں \textenglish{2021}، سوال \textenglish{6} اور سوال \textenglish{7})}}
اہلِ کتاب کی کتاب میں تحریف اور بخل کی مذمت میں آیات کا ترجمہ و تفسیر لکھیے۔
\end{sawal}

\begin{hal}
\ayat{وَإِنَّ مِنْهُمْ لَفَرِيقًا يَلْوُونَ أَلْسِنَتَهُم بِالْكِتَابِ
لِتَحْسَبُوهُ مِنَ الْكِتَابِ وَمَا هُوَ مِنَ الْكِتَابِ}{آل عمران،
\textenglish{78}}
\tarjuma{اور ان میں سے ایک گروہ کتاب پڑھتے وقت اپنی زبانیں مروڑتا ہے تاکہ تم
اسے کتاب کا حصہ سمجھو، حالانکہ وہ کتاب میں سے نہیں۔}

\smallskip
\textbf{تشریح:} یہ \emph{لفظی تحریف} کی تصویر ہے --- الفاظ کی ادائیگی بدل کر
یا من گھڑت عبارت ملا کر پڑھنا، تاکہ سننے والا اسے اصل کتاب کا حصہ سمجھ لے۔
یہی وہ عمل ہے جس نے تورات و انجیل کی اصل تعلیمات کو مسخ کیا۔

\medskip
\ayat{وَلَا يَحْسَبَنَّ الَّذِينَ يَبْخَلُونَ بِمَا آتَاهُمُ اللَّهُ مِن
فَضْلِهِ هُوَ خَيْرًا لَّهُم بَلْ هُوَ شَرٌّ لَّهُمْ سَيُطَوَّقُونَ مَا
بَخِلُوا بِهِ يَوْمَ الْقِيَامَةِ}{آل عمران، \textenglish{180}}
\tarjuma{اور جو لوگ اللہ کے دیے ہوئے فضل میں بخل کرتے ہیں، وہ ہرگز یہ گمان نہ
کریں کہ یہ ان کے لیے بہتر ہے، بلکہ یہ ان کے لیے بری ہے۔ عنقریب قیامت کے دن
انہیں اسی مال کا طوق پہنایا جائے گا جس میں انہوں نے بخل کیا۔}

\smallskip
\textbf{تشریح:} بخل یعنی اللہ کے دیے مال یا علم کو اس کے حق داروں سے روک
لینا۔ آیت میں \emph{طوق} کی تشبیہ خاص ہے --- جو مال دنیا میں دبا کر رکھا
گیا، وہی قیامت کے دن گلے کا طوق بن کر عذاب کا سبب بنے گا۔
\end{hal}

%---------------------------------------------------------------------
\begin{sawal}{سوال \textenglish{15} \ \textnormal{\small(خزاں \textenglish{2021}، سوال \textenglish{8})}}
کائنات کی تخلیق میں نشانیوں اور اولو الالباب (اہلِ عقل) کی صفت پر آیت کا
ترجمہ و تفسیر لکھیے۔
\end{sawal}

\begin{hal}
\ayat{إِنَّ فِي خَلْقِ السَّمَاوَاتِ وَالْأَرْضِ وَاخْتِلَافِ اللَّيْلِ
وَالنَّهَارِ لَآيَاتٍ لِّأُولِي الْأَلْبَابِ}{آل عمران، \textenglish{190}}
\tarjuma{بے شک آسمانوں اور زمین کی پیدائش میں، اور رات دن کے بدلتے رہنے میں،
عقل والوں کے لیے نشانیاں ہیں۔}

\smallskip
\ayat{الَّذِينَ يَذْكُرُونَ اللَّهَ قِيَامًا وَقُعُودًا وَعَلَىٰ جُنُوبِهِمْ
وَيَتَفَكَّرُونَ فِي خَلْقِ السَّمَاوَاتِ وَالْأَرْضِ رَبَّنَا مَا خَلَقْتَ
هَٰذَا بَاطِلًا سُبْحَانَكَ فَقِنَا عَذَابَ النَّارِ}{آل عمران،
\textenglish{191}}
\tarjuma{جو کھڑے، بیٹھے اور کروٹوں کے بل لیٹے اللہ کو یاد کرتے ہیں، اور
آسمانوں و زمین کی تخلیق میں غور کرتے ہیں: اے ہمارے رب! تو نے یہ سب بے مقصد
نہیں بنایا، تو پاک ہے، ہمیں آگ کے عذاب سے بچا۔}

\medskip
\textbf{تشریح:} \emph{اولو الالباب} یعنی وہ صاحبِ عقل جو کائنات کو محض دیکھتے
نہیں، اس پر \textbf{تفکر} (غور و فکر) بھی کرتے ہیں۔ ان کی تین صفات: ہر حال میں
(کھڑے، بیٹھے، لیٹے) اللہ کی یاد، کائنات میں غور، اور اس غور کا نتیجہ ---
اعترافِ عظمتِ الٰہی اور دعا۔ یہ آیات علم اور عبادت کے امتزاج کی بہترین مثال
ہیں: مشاہدہ، تفکر اور تسبیح ایک ساتھ۔
\end{hal}